\section{Представление деревьев в ЭВМ}
\label{sec:tree-representation}

\subsection{Почему дерево удобно хранить отдельно от общего графа}

Дерево можно представить списками смежности как обычный граф, однако корневая
иерархия даёт дополнительную информацию: у каждой вершины, кроме корня, ровно
один родитель. Поэтому многие структуры хранения имеют линейный объём
$O(n)$ и дают быстрые переходы вверх или вниз по иерархии.

При выборе представления важны операции:
\begin{itemize}
 \item переход к родителю и детям;
 \item обход поддерева;
 \item вставка/удаление вершины или поддерева;
 \item вычисление глубины и предков;
 \item поиск наименьшего общего предка (LCA);
 \item сериализация и восстановление дерева.
\end{itemize}

\subsection{Массив родителей}

Для вершин $0,\ldots,n-1$ хранится
\[
    parent[v].
\]
Для корня используют $-1$ или специальное значение. Память $O(n)$, переход к
родителю $O(1)$. Но без дополнительного индекса поиск всех детей требует полного
просмотра массива, то есть $O(n)$.

Массив родителей особенно удобен для задач, где часто поднимаются к корню.
Например, высота вершины может быть найдена цепочкой
\[
    v,\ parent[v],\ parent[parent[v]],\ldots.
\]
В худшем случае это $O(n)$, поэтому для многих запросов строят дополнительную
структуру.

\subsection{Двоичный подъём}

Пусть
\[
    up[v][k]
\]
--- предок вершины $v$ на расстоянии $2^k$. Тогда
\[
    up[v][0]=parent[v],
    \qquad
    up[v][k]=up[up[v][k-1]][k-1].
\]
Таблица строится за
\[
    O(n\log n)
\]
и занимает столько же памяти. Подъём на $d$ уровней выполняется по двоичному
разложению $d$ за $O(\log n)$.

Для LCA двух вершин сначала более глубокую вершину поднимают до уровня второй,
а затем одновременно пытаются поднимать обе вершины степенями двойки, начиная с
наибольшей. Запрос LCA занимает $O(\log n)$.

\subsection{Списки детей}

Для каждой вершины хранится динамический массив/список её детей. Так как в
дереве $n-1$ ребро, суммарная длина всех списков равна $n-1$, а память ---
$O(n)$.

Это представление идеально для обхода сверху вниз и обработки поддеревьев.
Если дерево статично, можно использовать компактный CSR-подобный формат:
\[
 children[offset[v]\ldots offset[v+1]).
\]
Он избегает большого числа маленьких динамических объектов и хорошо использует
кэш.

\subsection{Представление ``первый ребёнок --- следующий брат''}

Для дерева произвольной степени можно хранить в каждом узле только две ссылки:
\[
    firstChild(v),\qquad nextSibling(v).
\]
Первая ведёт к первому ребёнку, вторая --- к следующему ребёнку того же родителя.
Таким образом, произвольное корневое дерево кодируется структурой с постоянным
числом ссылок на вершину, похожей на бинарное дерево.

Преимущество --- компактная универсальная структура. Недостаток --- поиск $k$-го
ребёнка требует последовательного прохода по братьям.

\subsection{Связное представление бинарного дерева}

Классический узел:
\begin{verbatim}
struct Node {
    Data value;
    Node* left;
    Node* right;
    Node* parent; // optional
};
\end{verbatim}

Переход к левому/правому ребёнку $O(1)$. Указатель на родителя не обязателен, но
удобен для подъёма. Память $O(n)$, однако есть накладные расходы на указатели и
аллокации.

Для динамических деревьев это представление гибко: легко вставлять и отсоединять
поддеревья, если корректно поддерживаются ссылки.

\subsection{Последовательное представление полного/почти полного бинарного дерева}

Для деревьев, близких к полным, особенно для бинарных куч, узлы хранятся в
массиве. При нумерации с нуля:
\[
 left(i)=2i+1,
 \qquad
 right(i)=2i+2,
 \qquad
 parent(i)=\left\lfloor\frac{i-1}{2}\right\rfloor.
\]
При нумерации с единицы:
\[
 left(i)=2i,
 \qquad right(i)=2i+1,
 \qquad parent(i)=\left\lfloor\frac i2\right\rfloor.
\]

Плюсы: нет указателей, данные лежат компактно, отличная локальность памяти.
Минус: для сильно разреженного бинарного дерева прямое размещение по индексам
создаёт большое число пустых позиций.

\subsection{Обходы дерева}

Для бинарного дерева:
\begin{itemize}
 \item preorder: корень, левое, правое;
 \item inorder: левое, корень, правое;
 \item postorder: левое, правое, корень;
 \item level-order: по уровням с помощью очереди.
\end{itemize}
Каждый полный обход занимает $O(n)$ времени.

Если все ключи различны, пары
\[
    preorder+inorder
\]
или
\[
    postorder+inorder
\]
однозначно определяют бинарное дерево. Один preorder сам по себе недостаточен,
но preorder вместе с маркерами пустых детей позволяет однозначную сериализацию.

\subsection{Эйлеров обход корневого дерева}

Во время DFS запишем время входа $tin[v]$ и выхода $tout[v]$. Тогда вершина $u$
является предком $v$ тогда и только тогда, когда
\[
    tin[u]\le tin[v]
    \quad\text{и}\quad
    tout[v]\le tout[u].
\]

Если при входе в вершину присваивать ей номер $tin[v]$, вершины любого поддерева
образуют непрерывный диапазон в порядке DFS. Это позволяет сводить запросы по
поддереву к запросам на отрезке массива и использовать Fenwick tree, segment tree
и другие структуры данных.

Полный Euler tour, где вершина записывается при каждом посещении, также
используется для сведения LCA к задаче RMQ по глубинам.

\subsection{Код Прюфера}

Свободное помеченное дерево на $n$ вершинах кодируется последовательностью длины
$n-2$. Повторять:
\begin{enumerate}
 \item выбрать лист с минимальной меткой;
 \item записать метку его соседа;
 \item удалить лист.
\end{enumerate}

Код однозначен. Частота метки $v$ равна
\[
    d(v)-1.
\]
Это удобно для компактной передачи помеченного дерева и даёт формулу Кэли.
Кодирование/декодирование с очередью с приоритетом занимает $O(n\log n)$;
существуют линейные реализации при более специализированной организации данных.

\subsection{Деревья выражений}

В дереве арифметического выражения внутренние вершины --- операции, листья ---
операнды. Preorder даёт префиксную запись, postorder --- постфиксную. Например,
\[
    (a+b)c
\]
имеет постфиксную форму
\[
    a\ b\ +\ c\ *.
\]
Постфиксную запись можно вычислять стеком без скобок; такие деревья используются
в компиляторах и системах символьной алгебры.

\subsection{Что важно сказать на экзамене}

Основные представления: массив родителей, списки детей, ``первый ребёнок ---
следующий брат'', указатели left/right и массив для полного бинарного дерева.
Нужно сравнить память и стоимость переходов. В расширенном ответе хорошо показать
двоичный подъём/LCA, DFS-нумерацию поддеревьев и сериализацию обходами или кодом
Прюфера.

\newpage
